\section{Results}\label{sec:results}

\subsection{Raw differences}

Table \ref{tab:sumstats} shows why raw comparisons mislead. In La Liga, Barcelona receives 1.86 yellow cards per match against 2.53 for the other La Liga clubs, commits 11.6 fouls against 14.8, and suffers 15.3 fouls against 14.5. It is also awarded 0.19 penalties per match against 0.15 and concedes 0.10 against 0.16. Every one of these gaps favors Barcelona. But Barcelona also holds 66.8 percent of possession, takes 15.8 shots per match, enters matches with a 65 percent implied win probability, and wins by 1.46 goals on average. Real Madrid (column 2) looks similar on most margins; it is awarded more penalties than Barcelona (0.23 per match). Barcelona's per-match decisions in the Champions League (column 6) are strikingly close to its La Liga numbers: 1.71 yellow cards received, 15.3 fouls suffered, 0.20 penalties awarded. This is the first indication that much of the raw gap travels with the club rather than with Spanish referees.

\subsection{Conditional gaps}

Table \ref{tab:teamfe} reports estimates of equation (\ref{eq:teamfe}). In Panel A, conditional on home status, opponent, league-season and pre-match win probability, Barcelona receives 0.23 fewer yellow cards per match than the average non-elite club (column 1, standard error 0.06). It has 1.73 fewer fouls called against it (column 5) and 0.88 more fouls called on its opponents (column 6). There is no difference in yellow cards shown to its opponents (column 2), red cards (columns 3--4), or penalties awarded or conceded (columns 7--8; point estimates $-0.016$ and $-0.008$). Columns (9) and (10) are the most direct measure of referee discretion. Conditional on the number of fouls called, Barcelona's players are no less likely to be booked ($-0.056$, s.e. 0.054), and its opponents are no more likely to be booked ($-0.063$, s.e. 0.074). In other words, the yellow-card gap in column (1) is fully accounted for by the smaller number of fouls called against Barcelona.

Panel B adds possession-bin fixed effects and shots. The yellow-card gap falls to $-0.17$ and the fouls-against gap to $-1.28$; the fouls-for gap is unchanged at 0.94. Real Madrid's fouls-for gap (0.93) is virtually identical to Barcelona's, while its fouls-against gap is smaller ($-0.44$; $p = 0.016$ for equality with Barcelona).

Table \ref{tab:exposure} addresses the remaining foul gap by normalizing fouls by exposure. When fouls are divided by possession minutes (columns 1--2), the sign of Barcelona's advantage reverses. Barcelona suffers 0.17 \emph{fewer} fouls per ten minutes of possession and commits 0.25 \emph{more} fouls per ten minutes of opponent possession than the average club. The proportional-exposure Poisson models (columns 3--4) give the same result: 4 percent fewer fouls suffered, 5 percent more committed. Imposing proportionality is too strong, however. When the elasticity is estimated (columns 5--6), fouls rise by only 0.17 percent for each 1 percent of possession. At that elasticity, Barcelona suffers 4.6 percent more fouls and commits 9.4 percent fewer than predicted. Normalizing by passes (columns 7--8) yields no difference at all. Barcelona's foul advantage is therefore modest, at most about 5--10 percent, and its sign depends on how exposure is measured.

\subsection{Barcelona in the cross-club distribution}

Figure \ref{fig:ranks} places Barcelona among all 144 clubs with at least 150 domestic matches. On net yellow cards, Barcelona ranks 11th (13th with style controls), so about 8--9 percent of clubs are at least as favored. On net fouls it ranks 2nd (4th). On net red cards it ranks 56th (125th), and on net penalties 93rd of 150 (31st of 139). The clubs that share the top of the yellow-card and foul rankings are telling: Swansea City, Real Sociedad, Las Palmas, Girona, Celta Vigo, Lorient, Guingamp, Sassuolo. These are technically oriented, possession-minded clubs, mostly small, none of which is a plausible beneficiary of systematic favoritism. Real Madrid ranks 44th on net yellow cards and 11th on net fouls. The other elite clubs are spread throughout the distribution, with Paris Saint-Germain and Bayern Munich near the bottom on yellow cards. Being a large, successful club does not by itself buy favorable cards. On the most consequential decision, penalties, Barcelona is unremarkable.

\subsection{Barcelona against Europe's elite clubs, season by season}

Figures \ref{fig:elite-yellow}--\ref{fig:elite-heatmap} extend equation (\ref{eq:eventstudy}) to all five leagues. They estimate a separate gap for each of the fourteen elite clubs in each season, relative to the average non-elite club in that club's own league-season. Before 2018/19, Barcelona is the most favored elite club on both card and foul margins. Its net yellow-card gap averages 0.46 per match, against 0.24 for the next club (Borussia Dortmund) and 0.12 for Real Madrid. Its net foul gap averages 2.71, against 0.97 for the next club (Bayern Munich) and 0.78 for Real Madrid. After 2017/18, Barcelona's net yellow-card gap ($-0.24$) is in the middle of the elite distribution. Its net foul gap (1.27) is second only to Real Madrid's, which rises to 2.44. On penalties, no elite club, Barcelona included, shows a gap distinguishable from the season-to-season noise. The two Spanish giants are the only elite clubs with large and persistent positive foul margins; Atl\'etico Madrid's is close to zero. This suggests that the foul margin partly reflects how La Liga opponents defend against dominant sides, not only how referees treat one club.

\subsection{Spanish versus UEFA referees, before and after 2018}

Table \ref{tab:ddd} compares Barcelona's domestic decisions with its Champions League decisions, relative to the same gap for the other Spanish Champions League regulars. Before 2018/19, Barcelona received 0.38 fewer yellow cards per match in La Liga than its European baseline implies (Panel A, column 1; 0.50 with possession bins). It also had 1.1 fewer fouls called against it (column 3). After 2017/18, both gaps shrink toward zero: the yellow-card gap rises by 0.43 (s.e. 0.26) in Panel A and by 0.61 (s.e. 0.27) in Panel B. Two features temper a favoritism reading. First, Barcelona's \emph{opponents} also received fewer yellow cards in La Liga in the pre period ($-0.20$ and $-0.29$, column 2). The pattern is closer to lenient refereeing of Barcelona's matches than to one-sided treatment. Second, penalties, the decision with the largest effect on results, show no domestic premium in either period (columns 5--6). Real Madrid shows the opposite time pattern: its opponents' cards and fouls rise after 2018 (columns 2 and 4).

Figure \ref{fig:es} traces the timing in La Liga. Barcelona's net yellow-card advantage over the average club averages 0.36 per match through 2017/18 and falls by 0.68 (s.e. 0.13) afterward. Real Madrid's changes by only $-0.11$, so the difference-in-differences is $-0.57$ (s.e. 0.18). The decline is concentrated in 2019/20--2023/24, the seasons of Barcelona's sporting and financial crisis, and reverses in 2024/25--2025/26. The net-foul series does not break in 2018 at all. It stays near 2.5 per match through 2020/21 and collapses in 2021/22, the first season after Lionel Messi's departure. Penalties show no break (difference-in-differences $-0.04$, s.e. 0.07). The timing therefore fits changes in Barcelona's squad and style at least as well as the end of the payments. VAR, introduced in La Liga in 2018/19, cannot review yellow cards and so cannot directly explain the yellow-card change. It can, however, affect penalties, which show no break.

\subsection{Stoppage time}

Table \ref{tab:stoppage} reports the stoppage-time test. Across the five leagues, referees add 0.20 minutes more when the home team trails by one than when it leads by one, in line with \citet{garicano2005favoritism}. For Barcelona, the corresponding difference is 0.71 minutes (s.e. 0.24), roughly three and a half times the home-team effect. Real Madrid's is the same, 0.70 minutes (s.e. 0.22). Within La Liga, Barcelona's stoppage-time advantage is small and insignificant in 2005/06--2017/18 (0.20 minutes, s.e. 0.20). It is larger after 2018/19 (1.06 minutes, s.e. 0.45), again matched by Real Madrid (1.18). This is the clearest evidence of favoritism in the paper. It is shared by Spain's two dominant clubs and does not track the Negreira window, which points to the status-based social pressure documented by \citet{garicano2005favoritism} rather than to a Barcelona-specific arrangement.

\subsection{Real Madrid in the Champions League}

Real Madrid won the Champions League seven times in the sample period (2001/02, 2013/14, 2015/16, 2016/17, 2017/18, 2021/22, 2023/24), which makes it a natural case for testing whether UEFA referees favor the competition's dominant club. Table \ref{tab:rmucl} reports its decisions season by season. Over all seasons, Real Madrid receives 1.85 yellow cards per match and its opponents 1.96. It commits 11.6 fouls and suffers 13.4. It is awarded 0.17 penalties per match and concedes 0.13. These rates are close to those of the other elite clubs (1.76 and 1.89 yellow cards; 0.19 and 0.13 penalties). Barcelona's opponents are booked more often (2.42 per match).

Table \ref{tab:rmuclreg} reports conditional gaps from equation (\ref{eq:rmucl}). Relative to the average non-elite team in the same season and stage, facing the same opponent at the same Elo expected score and possession, Real Madrid receives 0.14 \emph{more} yellow cards per match (s.e. 0.10). This is significantly more than the other elite clubs ($p = 0.005$). Its opponents receive 0.10 fewer (s.e. 0.10). Real Madrid has 0.88 fewer fouls called against it (s.e. 0.44), but no more fouls called on its opponents, and no difference in penalties awarded ($-0.02$) or conceded ($-0.02$). In the knockout rounds, where titles are decided, Real Madrid receives 0.21 more yellow cards than comparable teams (s.e. 0.12), and its penalty gaps are small and negative. Across eras, the only large gap is in fouls called against Real Madrid in 2013/14--2017/18 ($-2.52$, s.e. 0.54). That gap is not accompanied by fewer yellow cards (0.02) or more penalties (0.02). Since 2018/19, Real Madrid's opponents have received 0.36 \emph{fewer} yellow cards per match (s.e. 0.11), a gap that runs against the team.

Figure \ref{fig:rmucl} plots the season-by-season gaps. Title seasons are not systematically favorable. Real Madrid's net yellow-card gap is statistically indistinguishable from zero in all six title seasons with card data (estimates between $-0.35$ and 0.57). The net foul margin is significantly positive in two title seasons, 2016/17 (5.4 per match) and 2021/22 (2.9). Penalties cut both ways: the net penalty gap is $-0.41$ per match in 2016/17 (five conceded, none awarded) and $+0.28$ in 2017/18.

Table \ref{tab:rmuclstoppage} applies the stoppage-time test to close Champions League matches. Across all stages, referees add 0.55 more minutes when Real Madrid trails by one at 90 minutes than when it leads (s.e. 0.39). The effect is concentrated in the group/league phase (1.47 minutes, s.e. 0.48). In the knockout rounds, it is $-0.86$ (s.e. 0.63), based on 20 matches in which Real Madrid led and 14 in which it trailed. Barcelona's corresponding gap is 0.94 minutes (s.e. 0.42) and the other elite clubs' is 0.29 (s.e. 0.12). The pattern resembles the domestic results: big clubs receive more time when trailing, with no special advantage for Real Madrid where the stakes are highest.
